\chapter{Wet en Regelgeving}

't Jansje krijgt een aandrijving waarvan niet alle onderdelen een nautische goedkeuring hebben. Dit betekent dat deze aandrijving niet zonder meer in een schip gebruikt mag worden. Gelukkig zijn er wel manieren om een goedkeuring te krijgen in specifieke gevallen. In dit hoofdstuk zal onderzocht worden hoe 't Jansje aanspraak kan maken op deze goedkeuring.

\subsection{Wet \& regelgeving experimentele batterijen}
Om te beginnen moet er normaalgesproken een algemene typegoedkeuring worden gegeven voor het schip. Echter verschilt deze per zone van de wereld. Uit een interview met de eigenaar van het schip wordt impliciet de conclusie getrokken dat het schip niet buiten Nederlandse binnenlandse wateren zal varen. Nederland valt onder Zone 2 \cite{ministerie2026inlandrequirements}. 


\subsection{algemene Wet \& regelgeving voor elektrisch aangedreven schepen}

Alle volgende informatie komt van CESNI (Europees comité voor het opstellen van normen voor de binnenvaart)\cite{cesni2025inlandstandards}.
Allereerst wordt er gekeken naar de eisen voor motorontwerp van schepen.
Hierin staan de volgende relevante eisen:
\begin{itemize}
    \item motoren zullen geplaatst worden in een dergelijke manier dat het opereren dan wel onderhouden van het schip geen gevaar brengt aan de personen die deze taken uitvoeren (ES-TRIN 2025 8.02 1).
    \item het moet mogelijk zijn om het voortstuwingssysteem van het schip uit te schakelen van buiten de motorruimte (ES-TRIN 2025 8.02 3).
    \item Bij een schip met één motor, mag een automatische uitschakeling niet mogelijk zijn, tenzij deze beschermt tegen overtoeren (ES-TRIN 2025 8.03 2).
    \item asbussen moeten van dergelijke ontwerp zijn dat het verspreiden van watervervuilende glijmiddelen wordt verkomen (ES-TRIN 2025 8.03 4).\\\\
    \item Voor elektrsiche aandrijvingssystemen zal ten alle tijden een process- en instrumentatiediagram (PID) aanwezig zijn, welke is goedgekeurd door de instantie van inspectie (ES-TRIN 2025 10.01 2e)
    \item elektrische systemen zullen geplaatst worden in een dergelijke manier dat het opereren dan wel onderhouden van het schip geen gevaar brengt aan de personen die deze taken uitvoeren (ES-TRIN 2025 10.01 4).
    \item er moeten ten minste twee elektrische voedingssystemen aanwezig zijn op het schip zodat, mocht één voeding uitvallen, de tweede nog steeds veilige operatie van het ship garandeert voor ten minste 30 minuten (ES-TRIN 2025 10.02 1).
    \item de adequate dimensionering van de stroomvoorziening  moet worden aangetoond door middel van een energiebudgetberekening (ES-TRIN 2025 10.02 2)
    \item batterijen met een laadvermogen boven 2kW moeten geïnstalleerd worden in een speciale ruimte of op het dek (ES-TRIN 2025 10.11 3).
    \item het laadvermogen is berekend gebaseerd op de maximale laadstroom en de nominale spanning van de batterij (Es-TRIN 2025 10.11 4).
    \item alle materialen in deze ruimte moeten beschermd zijn tegen beschadigende effecte van het elektrolyt in de batterij (ES-TRIN 2025 10.11 6).
    \item er zullen waarschuwingen voor 'vuur', 'open vlam' en 'roken verboden' op alle deuren naar de batterijkamer kenbaar gemaakt worden.
    \item er moeten vuurwerende, zelfdovende elektriciteitskabels gebruikt worden die bestand zijn tegen olie en water (ES-TRIN 2025 10.15 1).
\end{itemize}

Verder is voor de zekerheid aan te raden om de ruimte met batterijen mechanische te ventileren, vergelijkbaar met ES-TRIN 2025 10.11 7.

Ook moeten er een aantal tests uitgevoerd worden op de elektronica:
\begin{itemize}
    \item Variatie in stroom en spanning (ES-TRIN 2025 10.20 2b)
    \item hittetest (ES-TRIN 2025 10.20 2c)
    \item koudetest (ES-TRIN 2025 10.20 2d)
    \item vibratietest (ES-TRIN 2025 10.20 2e37
    )
\end{itemize}

Uiteindelijk is de conclusie dat het implementeren van deze regelgeving belangrijk is, maar buiten de omvang van dit project valt. Het toepassen van deze regelgeving zal een aanbeveling zijn voor vervolgstappen.
